% =====================================================================
\chapter{Tổng kết}\label{chap:ketluan}
% =====================================================================

\section{Kết quả đạt được}
Nhóm đã xây dựng hoàn chỉnh một hệ thống bãi đỗ xe thông minh quy mô mô hình gồm ba
vi điều khiển, cùng kết nối vào mạng Wi-Fi chung của nhóm:
\begin{itemize}
  \item \textbf{10 use case} (Chương~\ref{chap:usecase}) đều được hiện thực: theo dõi chỗ
        trống tại bãi và trên web, gửi/lấy xe bằng RFID kèm tính phí, từ chối quẹt sai
        cổng hoặc bãi đầy, cảnh báo cháy và mở đường sơ tán, giám sát nước và nhiệt độ,
        điều khiển thanh chắn thủ công, nhật ký vào/ra kèm phí.
  \item \textbf{Thanh toán VietQR}: cổng ra hiển thị mã VietQR (BIDV) sinh ngay trên chip,
        máy chủ hỏi SePay mỗi 3\,s và tự mở thanh chắn khi nhận được chuyển khoản đúng mã
        và đủ tiền; khi không có Internet hoặc trả tiền mặt, người vận hành xác nhận trên
        trang web.
  \item \textbf{Kiến trúc} phân lớp rõ ràng, không có phụ thuộc vòng; đa tác vụ với
        FreeRTOS và ba cơ chế đồng bộ dùng đúng chỗ; giao tiếp UDP nhẹ, tự phục hồi khi mất gói;
        tên \code{parking.local} luôn trỏ đúng IP; ba bo đều là trạm Wi-Fi nên bật theo
        thứ tự nào cũng được.
  \item \textbf{Hiệu năng}: màn hình phản ánh thay đổi trong khoảng 0,1\,s; thanh chắn mở
        khoảng 0,2\,s sau khi quẹt thẻ; firmware dùng 31\% RAM tĩnh, còn 20\% flash ứng dụng.
  \item \textbf{Chất lượng}: 384 phép kiểm thử tự động đều đạt, chạy cả mã nguồn hai
        client trên phần cứng giả lập.
  \item \textbf{Bền vững}: quẹt thẻ không bị mất (ACK và gửi lại), dữ liệu gửi xe giữ được qua mất điện, giờ thực không cần Internet.
  \item \textbf{Giám sát chủ động}: phát hiện Sensor Client mất kết nối sau 3\,s thay vì
        hiển thị dữ liệu cũ như dữ liệu thật; cảnh báo ngập có trễ, tự chặn xe vào; thống kê
        doanh thu, lượt vào/ra theo ngày và 7 ngày.
  \item \textbf{Mô hình thực tế} hoàn chỉnh (Hình~\ref{fig:model}): 5 ô có cảm biến, 2 cổng
        servo, hộp điều khiển trung tâm có màn hình, nguồn có bảo vệ ngược cực.
\end{itemize}

Qua đề tài, nhóm đã vận dụng các kiến thức trọng tâm của môn lập trình nhúng: hệ điều
hành thời gian thực, đồng bộ tác vụ, PWM và SPI có DMA, ADC và lọc nhiễu tín hiệu số,
lập trình socket, quản lý bộ nhớ hạn chế, xử lý dữ liệu không tin cậy, và kiểm thử
phần mềm nhúng không cần phần cứng.

\section{Các hạn chế đã được khắc phục}
Bản thiết kế đầu tiên có năm hạn chế. Bốn hạn chế đã được khắc phục hoàn toàn, hạn
chế về bảo mật được khắc phục ở phía trang web và mật khẩu (bản tin UDP chủ đích không
ký, xem mục~\ref{sec:conlai}); tất cả được
kiểm chứng bằng kiểm thử tự động (Bảng~\ref{tab:unittests}) và kịch bản phần cứng
S15–S21 (Bảng~\ref{tab:systest}):

\begin{longtable}{@{}p{3.2cm}p{6.6cm}p{4.6cm}@{}}
\caption{Hạn chế ban đầu và giải pháp đã hiện thực}\label{tab:fixed}\\
\toprule \textbf{Hạn chế} & \textbf{Giải pháp} & \textbf{Kiểm chứng} \\ \midrule \endfirsthead
\toprule \textbf{Hạn chế} & \textbf{Giải pháp} & \textbf{Kiểm chứng} \\ \midrule \endhead
Bản tin UDP và API không có xác thực; mật khẩu cố định trong mã nguồn &
Lệnh POST cần mã PIN quản trị; mật khẩu Wi-Fi và mã PIN chuyển ra \file{sdkconfig}
(menuconfig) và \file{secrets.h}, không còn trong Git &
S19 \\ \midrule
Sự kiện quẹt thẻ có thể mất &
Máy chủ trả ACK; RFID Client gửi lại sau 300\,ms, tối đa 6 lần; máy chủ nhận ra
bản gửi lại qua \code{seq} nên không mở cổng hai lần &
\file{rfid\_client}: gửi lại, bỏ cuộc; \file{rfid\_session}: chống xử lý trùng; S16, S18 \\ \midrule
Phiên và nhật ký mất khi khởi động lại; thời gian tính từ lúc khởi động &
Toàn bộ trạng thái (\code{rfid\_store\_t}) lưu NVS sau mỗi lần quẹt và mỗi 5 phút, khôi
phục khi khởi động; trình duyệt của người vận hành gửi giờ Unix
(\code{POST /api/time}), máy chủ dời các mốc cũ theo và bù thời gian mất điện &
\file{rfid\_session}: lưu/khôi phục, đặt giờ; S17 \\ \midrule
Nhiệt độ chỉ cập nhật khi có xe quẹt thẻ &
RFID Client gửi bản tin \code{climate} mỗi 5\,s, độc lập với quẹt thẻ &
\file{rfid\_client}, \file{json\_parsers}; S21 \\ \midrule
Số phiên RFID và cảm biến IR chưa được đối chiếu &
Máy chủ từ chối thẻ mới khi IR báo đủ 5 ô có xe; trang web cảnh báo \code{CHECK THE LOT}
khi số thẻ trong bãi lệch số ô có xe quá 15\,s &
\file{rfid\_session}: đối chiếu IR; S20 \\
\bottomrule
\end{longtable}

\section{Hạn chế còn lại}\label{sec:conlai}
Các điểm dưới đây là đánh đổi có chủ đích ở quy mô mô hình, đã được ghi rõ trong mã
nguồn và hướng dẫn:
\begin{itemize}
  \item Trang web dùng HTTP (không mã hóa) trên mạng WPA2: người trong cùng mạng Wi-Fi có
        thể nghe được mã PIN. HTTPS cần chứng chỉ và thêm $\approx$ 40\,KB RAM cho TLS.
  \item Bản tin UDP không được ký: thiết bị đã vào được mạng Wi-Fi có thể gửi bản tin giả
        (thẻ giả, báo cháy giả). Nhóm chấp nhận đánh đổi này vì mạng là mạng Wi-Fi riêng của
        nhóm có WPA2 và giao thức đơn giản dễ gỡ lỗi hơn.
  \item Bãi phụ thuộc vào router/điểm phát Wi-Fi chung: router hỏng thì các bo không nói
        chuyện được với nhau; router bật chế độ cô lập thiết bị sẽ chặn gói broadcast.
  \item Token SePay có toàn quyền đọc giao dịch của tài khoản (SePay không giới hạn quyền
        token) và nằm trong flash của máy chủ: ai đọc được flash là lấy được token.
        Mỗi lúc chỉ một xe được thanh toán ở cổng ra.
  \item Nếu chưa có trình duyệt nào gửi giờ sau khi mất điện, thời gian mất điện chưa
        được tính vào thời gian gửi xe (được bù ngay khi trình duyệt kết nối).
\end{itemize}

\section{Hướng phát triển}
\begin{itemize}
  \item Cảm biến IR dưới mỗi thanh chắn: chỉ đóng khi làn xe trống, tránh hạ cần xuống xe.
  \item Danh sách thẻ đăng ký (vé tháng, khóa thẻ mất) lưu trong NVS, quản lý từ trang web.
  \item Ghép thẻ với ô đỗ: ô đầu tiên chuyển sang có xe sau khi quẹt vào thuộc về thẻ đó.
  \item Ký bản tin UDP bằng HMAC kèm số thứ tự chống phát lại; HTTPS cho trang quản trị.
  \item Module RTC DS3231 để có giờ thực ngay cả khi chưa có trình duyệt kết nối.
  \item Kết nối MQTT lên đám mây để quản lý nhiều bãi, thống kê doanh thu, đặt chỗ trước.
  \item Bổ sung OTA cho máy chủ (bảng phân vùng hai khe ứng dụng, cần flash 4\,MB).
\end{itemize}
